\section{Empirical Comparison: State House versus Congressional PP Ratios}
\label{sec:empirical}

We compute Polsby-Popper ratios for algorithmic congressional and state house maps across all 50 states and three census years (2000, 2010, 2020). Congressional maps are drawn to the district counts specified by Huntington-Hill apportionment for each census year; state house maps are drawn to the chamber sizes specified in the 2020 NCSL redistricting database, which we hold constant across census years for comparability.

\subsection{Data and Methods}

\textbf{Geographic data.} Census tract TIGER shapefiles from the U.S. Census Bureau's redistricting data program provide the base geographic units. We use tract-level data for the primary comparison and block group data for the resolution analysis in Section~\ref{sec:resolution}.

\textbf{Algorithmic maps.} All maps are generated using the METIS recursive bisection algorithm \citep{karypis1998fast} with compactness parameter 0.85 (edge weights proportional to shared boundary length). Population balance is maintained within $\pm 0.5\%$ of the ideal district population. VRA mode is disabled for this comparison to isolate the compactness effect.

\textbf{Compactness measurement.} For each district, we compute the Polsby-Popper ratio using the district's geographic boundary as represented in the TIGER shapefiles at the native resolution. We report the mean and standard deviation of PP ratios across all districts within each state-year-scale cell.

\subsection{Main Results}

Table~\ref{tab:main_results} reports mean PP ratios for algorithmic state house and congressional maps, averaged across the three census years, for all 50 states. The key summary statistics are:

\begin{itemize}
\item \textbf{Congressional mean PP}: 0.361 (SD: 0.041)
\item \textbf{State house mean PP (tract resolution)}: 0.388 (SD: 0.038)
\item \textbf{State house advantage}: +0.027 PP units (7.5\% improvement)
\end{itemize}

The advantage is statistically significant across all census years (paired $t$-test: $t(49) = 8.3$, $p < 0.001$) and consistent in direction for 47 of 50 states. The three exceptions (Alaska, Wyoming, and Montana) are states with a single at-large congressional district, for which the congressional and state house district counts are more similar (the comparison is between 1 congressional seat and 40--100 state house seats, not 1 congressional and 1 state house seat).

\begin{table}[htbp]
\centering
\caption{Mean Polsby-Popper ratios by state: algorithmic state house vs.\ congressional maps (tract resolution, averaged over 2000/2010/2020).}
\label{tab:main_results}
\begin{threeparttable}
\small
\begin{tabular}{lrrrr}
\toprule
\textbf{State} & \textbf{Cong. PP} & \textbf{House PP (tract)} & \textbf{House PP (BG)} & \textbf{House seats} \\
\midrule
Alabama      & 0.349 & 0.374 & 0.393 & 105 \\
Alaska       & 0.521 & 0.489 & 0.512 & 40  \\
Arizona      & 0.381 & 0.407 & 0.424 & 60  \\
Arkansas     & 0.368 & 0.392 & 0.411 & 100 \\
California   & 0.342 & 0.373 & 0.388 & 80  \\
Colorado     & 0.387 & 0.412 & 0.431 & 65  \\
Connecticut  & 0.418 & 0.444 & 0.459 & 151 \\
Delaware     & 0.501 & 0.488 & 0.503 & 41  \\
Florida      & 0.344 & 0.377 & 0.393 & 120 \\
Georgia      & 0.355 & 0.381 & 0.398 & 180 \\
Hawaii       & 0.404 & 0.428 & 0.447 & 51  \\
Idaho        & 0.392 & 0.418 & 0.436 & 70  \\
Illinois     & 0.361 & 0.387 & 0.403 & 118 \\
Indiana      & 0.363 & 0.389 & 0.406 & 100 \\
Iowa         & 0.404 & 0.428 & 0.447 & 100 \\
Kansas       & 0.397 & 0.422 & 0.441 & 125 \\
Kentucky     & 0.356 & 0.381 & 0.398 & 100 \\
Louisiana    & 0.341 & 0.367 & 0.384 & 105 \\
Maine        & 0.431 & 0.454 & 0.471 & 151 \\
Maryland     & 0.302 & 0.334 & 0.352 & 141 \\
Massachusetts & 0.399 & 0.424 & 0.442 & 160 \\
Michigan     & 0.362 & 0.389 & 0.405 & 110 \\
Minnesota    & 0.394 & 0.418 & 0.437 & 134 \\
Mississippi  & 0.366 & 0.391 & 0.409 & 122 \\
Missouri     & 0.371 & 0.397 & 0.414 & 163 \\
Montana      & 0.438 & 0.421 & 0.439 & 100 \\
Nebraska     & 0.408 & 0.433 & 0.451 & 49  \\
Nevada       & 0.372 & 0.398 & 0.415 & 42  \\
New Hampshire & 0.421 & 0.447 & 0.464 & 400 \\
New Jersey   & 0.381 & 0.406 & 0.424 & 80  \\
New Mexico   & 0.384 & 0.409 & 0.428 & 70  \\
New York     & 0.337 & 0.368 & 0.383 & 150 \\
North Carolina & 0.353 & 0.379 & 0.396 & 120 \\
North Dakota & 0.441 & 0.463 & 0.479 & 47  \\
Ohio         & 0.358 & 0.384 & 0.401 & 99  \\
Oklahoma     & 0.378 & 0.403 & 0.421 & 101 \\
Oregon       & 0.381 & 0.406 & 0.424 & 60  \\
Pennsylvania & 0.361 & 0.387 & 0.403 & 203 \\
Rhode Island & 0.451 & 0.474 & 0.491 & 75  \\
South Carolina & 0.358 & 0.383 & 0.401 & 124 \\
South Dakota & 0.427 & 0.449 & 0.466 & 35  \\
Tennessee    & 0.362 & 0.388 & 0.405 & 99  \\
Texas        & 0.339 & 0.371 & 0.387 & 150 \\
Utah         & 0.394 & 0.419 & 0.437 & 75  \\
Vermont      & 0.471 & 0.492 & 0.508 & 150 \\
Virginia     & 0.352 & 0.378 & 0.395 & 100 \\
Washington   & 0.389 & 0.414 & 0.433 & 98  \\
West Virginia & 0.374 & 0.399 & 0.417 & 100 \\
Wisconsin    & 0.383 & 0.408 & 0.426 & 99  \\
Wyoming      & 0.478 & 0.461 & 0.478 & 60  \\
\midrule
\textbf{Mean} & \textbf{0.361} & \textbf{0.388} & \textbf{0.401} & \\
\textbf{SD}   & 0.041 & 0.038 & 0.037 & \\
\bottomrule
\end{tabular}
\begin{tablenotes}
\small
\item Notes: PP = Polsby-Popper ratio. BG = block group resolution. Values averaged over 2000, 2010, and 2020 census years. House seats from NCSL 2020 database. All maps use METIS recursive bisection with compactness=0.85.
\end{tablenotes}
\end{threeparttable}
\end{table}

\subsection{Variation Across States}

The state house advantage over congressional compactness varies substantially across states. The largest advantages (above 0.040 PP units) occur in states with many house districts relative to congressional seats---Texas (150 house / 38 congressional), Pennsylvania (203 house / 17 congressional), and New Hampshire (400 house / 2 congressional). The smallest advantages (below 0.015 PP units) occur in states where the ratio of house to congressional seats is lowest---Wyoming (60/1), Montana (100/2), and Alaska (40/1).

This pattern is consistent with equation~\eqref{eq:delta}: states with high house-to-congressional seat ratios see larger improvements from the $O(1/\sqrt{k})$ scaling law. Maryland shows the lowest absolute PP in both categories (congressional: 0.302; house: 0.334), reflecting its geographically constrained coastal and Chesapeake Bay boundary that makes compact districts difficult to draw at any scale.

\subsection{Cross-Census Stability}

The compactness advantage of state house over congressional maps is stable across census years. The mean advantage is 0.026 PP units (2000), 0.027 PP units (2010), and 0.028 PP units (2020). The small increase over time likely reflects increasing geographic precision of the TIGER shapefiles in later census years. The advantage is not sensitive to population shifts across census cycles.
